\chapter{Storage}
\section{Tổng quan về Cỗ máy Lưu trữ (Storage Engine)}

Cỗ máy lưu trữ của KBMS (phiên bản V3) chịu trách nhiệm quản lý việc ghi dữ liệu tri thức xuống đĩa cứng một cách bền vững và truy xuất chúng với hiệu năng cao nhất.

\subsection{Cấu trúc Trang (Page Structure)}

KBMS chia file dữ liệu thành các đơn vị cố định gọi là \textbf{Page} (Trang), mặc định là \textbf{16KB}. Việc sử dụng kích thước trang cố định giúp tối ưu hóa việc đọc/ghi theo từng block của ổ cứng (HDD/SSD).

\subsubsection{Slotted Page Algorithm}
Để lưu trữ các bản ghi (Tuples) có độ dài thay đổi, KBMS sử dụng thuật toán \textbf{Slotted Page}.
*   \textbf{Header:} Lưu trữ Metadata trang (24 Bytes): PageId, LSN, PrevPageId, NextPageId, FreeSpacePointer, TupleCount.
*   \textbf{Slot Array:} Một mảng các "khe cắm" phát triển từ đầu trang về phía sau. Mỗi khe chứa `offset` (vị trí) và `length` (độ dài) của bản ghi.
*   \textbf{Tuples Data:} Dữ liệu thực tế được đẩy từ cuối trang ngược về phía trước.
*   \textbf{Ưu điểm:} Tận dụng tối đa không gian trống ở giữa trang mà không cần biết trước số lượng bản ghi.

---

\subsection{Phân tích Độ phức tạp (Complexity Analysis)}

Hệ thống lưu trữ V3 được tối ưu hóa để xử lý hàng triệu bản ghi với độ trễ tối thiểu (Low-latency).

| Thành phần | Thuật toán | Độ phức tạp | Ghi chú |
| :--- | :--- | :--- | :--- |
| \textbf{B+ Tree Search} | Binary Search (Nodes) | $O(\log\_b n)$ | $b$ là bậc của cây (Branching Factor). |
| \textbf{B+ Tree Insert} | Split \& Rebalance | $O(\log\_b n)$ | Đảm bảo cây luôn cân bằng tuyệt đối. |
| \textbf{B+ Tree Delete} | Merge \& Rebalance | $O(\log\_b n)$ | Giải phóng không gian trang khi cần. |
| \textbf{Buffer Look-up} | Hash Table | $O(1)$ | Truy xuất trang trong RAM cực nhanh. |
| \textbf{WAL Write} | Sequential Log | $O(1)$ | Tốc độ ghi nhanh vì không cần Disk Seek. |

---

\subsection{Quản lý Bộ đệm (Buffer Pool Management)}

Để giảm thiểu I/O (đọc/ghi file) chậm chạp, KBMS duy trì một vùng nhớ RAM gọi là \textbf{Buffer Pool}.

\subsubsection{Thuật toán LRU (Least Recently Used)}
KBMS sử dụng chính sách thay thế LRU để quản lý các trang trong bộ đệm:
1.  \textbf{Pinning:} Khi một trang đang được sử dụng (để đọc hoặc ghi), nó được "Pin" (ghim). Trang đã bị ghim sẽ \textbf{không bao giờ} bị đẩy ra khỏi RAM.
2.  \textbf{Eviction (Trục xuất):} Khi Buffer Pool đầy và cần nạp trang mới, hệ thống sẽ chọn trang có `PinCount = 0` và thời gian truy cập xa nhất để loại bỏ.
3.  \textbf{Dirty Pages:} Nếu trang bị loại bỏ đã bị sửa đổi (`IsDirty = true`), hệ thống sẽ tự động thực hiện lệnh `Flush` để ghi nội dung xuống đĩa trước khi ghi đè trang mới vào frame đó.

---

\subsection{Hạn mức Kỹ thuật (Technical Limits)}

KBMS V3 được thiết kế để mở rộng (Scale-out) mạnh mẽ với các thông số định lượng sau:

| Thông số | Giá trị thực tế | Giới hạn lý thuyết |
| :--- | :--- | :--- |
| \textbf{Dung lượng CSDL} | \textbf{~35.25 TB} | $2^{31} \times 16,384$ Bytes |
| \textbf{Số lượng Trang (Pages)}| \textbf{2.1 Tỷ trang} | $2,147,483,647$ (4-byte ID) |
| \textbf{Kích thước Trang} | \textbf{16 KB} | Cố định theo kiến trúc V3. |
| \textbf{Bậc của cây (B+ Tree)} | \textbf{~512} | Dựa trên Key 32B + PageId 4B. |
| \textbf{Cache RAM tối đa} | Theo tệp `.ini` | Giới hạn bởi OS (64-bit). |

---

\subsection{Chiến lược Luồng I/O (I/O Strategy)}

Dữ liệu luân chuyển qua 3 trạng thái để đảm bảo cả tốc độ và độ an toàn:

\begin{figure}[ht!]
\centering
\includegraphics[width=0.8\textwidth]{assets/io\_strategy.png}
\caption{Sơ đồ chiến lược luồng I/O đa tầng của KBMS.}
\label{fig:7_1}
\end{figure}

<details>
<summary>Cấu trúc Mermaid (Source)</summary>

\begin{verbatim}
graph LR
    User[User/Reasoning] -- Write Fact --> BP[Buffer Pool: RAM]
    BP -- Log Record --> WAL[WAL: Sequential Log]
    BP -- Async Flush --> Disk[Data File: Random Access]
    WAL -- Integrity Check --> Disk
\end{verbatim}
</details>

---

\subsection{Danh mục Đặc tả Tầng Lưu trữ}

Để tìm hiểu sâu hơn về kiến trúc nhị phân và cơ chế vận hành, hãy tham khảo:

- [02-Cấu trúc Cây B+ Tree](./02-btree.md)
- [03-Cấu trúc Trang vật lý](./03-page-layout.md)
- [04-Cơ chế Độ bền \& Phục hồi](./04-recovery.md)
- [05-Quản lý Đĩa \& Mã hóa](./05-disk-management.md)

---

\subsection{Nhật ký ghi trước (Write-Ahead Logging - WAL)}

Tất cả các thay đổi dữ liệu đều được ghi vào file Log trước khi được áp dụng vào file dữ liệu chính.

\subsubsection{Sơ đồ Luồng ghi WAL}
\begin{figure}[ht!]
\centering
\includegraphics[width=0.8\textwidth]{assets/wal\_log\_flow.png}
\caption{Cơ chế ghi nhật ký trước (WAL) để đảm bảo tính bền vững ACID.}
\label{fig:7_2}
\end{figure}

*   \textbf{Mục tiêu:} Đảm bảo tính \textbf{Durability} (Độ bền vững). Nếu hệ thống bị sập nguồn đột ngột, KBMS sẽ đọc file log để tái tạo (Redo) lại các giao dịch chưa kịp ghi xuống đĩa.
*   \textbf{LSN (Log Sequence Number):} Mỗi bản ghi Log và mỗi trang dữ liệu đều có một số thứ tự LSN để đảm bảo tính tuần tự và kiểm tra phiên bản.

---

\subsection{Định dạng Tệp tin}

Hệ thống lưu trữ dưới các định dạng file tùy chỉnh:
*   `.kbf` (Knowledge Base File): Lưu trữ metadata và catalog của KB.
*   `.dat`: Lưu trữ dữ liệu thực tế của các Concept dưới dạng cây B+ Tree.
*   `.log`: Lưu trữ nhật ký WAL.

---

\subsection{Kế hoạch Triển khai (Implementation Strategy)}

Để hoàn thiện tầng giao tiếp đĩa cứng và đảm bảo tính ACID, tiến trình lập trình Backend (Architecture Coding Approach) cho Storage Engine được thực thi nghiêm ngặt qua 4 giai đoạn cốt lõi:

\subsubsection{Giai đoạn 1: Quản lý Vật lý & Cơ chế Mã hóa (DiskManager)}
Khởi tạo cấu hình truy xuất vật lý `FileStream.Seek` cho phép nhảy thẳng đến khối bộ nhớ theo kích thước khối cố định (16KB). Tích hợp module mã hóa tĩnh `AES-256`, thiết lập vector khóa IV để mã hóa/giải mã on-the-fly mỗi khi lấy dữ liệu ra khỏi bề mặt đĩa cứng. Thiết lập các chốt an toàn Disk Allocation (kỹ thuật Zero-fill) để chống phân mảnh cấp độ File System OS.

\subsubsection{Giai đoạn 2: Bộ đệm & An toàn I/O (BufferPool & WalManager)}
Viết lớp quản lý RAM nội vụ (`BufferManager`). Cấu hình thuật toán điều phối `LRU - Least Recently Used`, lập trình cơ chế "Ghim" trang (`PinCount`) và đánh dấu bị sửa (`IsDirty`). Đồng thời, lập trình kết nối hệ thống Write-Ahead Logging. Khi có giao dịch bẩn, ghi file `.klf` bằng phương thức quét tuần tự (Sequential I/O), đảm bảo Forward Recovery hoạt động hoàn hảo khi bị cúp điện. 

\subsubsection{Giai đoạn 3: Giải phẫu Chuỗi Nhị phân (Slotted Page)}
Tiến lên cấu trúc trang học thuật. Lập trình bộ nén nhị phân (Binary Serialization) để tạc khối 16,384 Bytes thành các `Page`. Thiết kế `Header` (24 Bytes Metadata), cấu hình mảng "Slot Array" (phát triển ngược chiều Payload) nhằm lưu thông tin khoảng cách (Offset) cho Tuple độ dài tùy biến. Chia nhánh kế thừa lớp trang thành `LeafPage` (Nút Lá) và `InternalPage` (Nút Trung Gian).

\subsubsection{Giai đoạn 4: Hạ tầng Đồ thị Thuần túy (B+ Tree Index)}
Kết nối hàng triệu khối nhị phân thành một đồ thị khổng lồ. Lập trình hệ thống B+ Tree Indexing $O(\log n)$. Đẩy dữ liệu gốc trực tiếp vào Cấp Lá (Clustered Index Pattern) để tiết kiệm thao tác truy xuất. Khởi chạy và kiểm nghiệm thuật toán Đệ quy tự tách nhánh (Split Node) khi trang dữ liệu báo trạng thái cạn không gian (Overflow).

\section{Cấu trúc Cây B+ Tree}

KBMS sử dụng cây B+ Tree làm cấu trúc dữ liệu chính để quản lý và truy xuất các Fact (bản ghi) trong mỗi Concept. B+ Tree là lựa chọn tối ưu cho việc truy vấn dữ liệu từ bộ nhớ ngoài - Disk [5].

\subsection{Cấu trúc Đa tầng (Tiered Structure)}

Cây B+ Tree trong KBMS V3 thường duy trì từ 3 đến 5 tầng tùy thuộc vào khối lượng tri thức:

\begin{figure}[ht!]
\centering
\includegraphics[width=0.8\textwidth]{assets/btree\_graph.png}
\caption{Sơ đồ cấu trúc đa tầng của cây B+ Tree trong việc quản lý Fact.}
\label{fig:7_3}
\end{figure}

<details>
<summary>Cấu trúc Mermaid (Source)</summary>

\begin{verbatim}
graph TD
    Root[Root Node: Page 1] -- Child 0 --> Int1[Internal Node: Page 10]
    Root -- Child 1 --> Int2[Internal Node: Page 20]
    
    Int1 -- Child 0 --> L1[Leaf Node: Page 101]
    Int1 -- Child 1 --> L2[Leaf Node: Page 102]
    
    Int2 -- Child 0 --> L3[Leaf Node: Page 201]
    Int2 -- Child 1 --> L4[Leaf Node: Page 202]
    
    L1 --- L2 --- L3 --- L4
    subgraph Facts
        L1
        L2
        L3
        L4
    end
\end{verbatim}
</details>

---

\subsection{Quy trình Vận hành}
Fact (Tuple).
*   \textbf{Liên kết:} Mỗi Node lá đều có con trỏ `NextPageId` trỏ đến Node lá tiếp theo, cho phép hệ thống thực hiện quét dữ liệu (Scan) theo phạm vi (Range Scan) cực nhanh mà không cần duyệt lại từ gốc.

---

\subsection{Các Thuật toán Cốt lõi}

\subsubsection{Thuật toán Tìm kiếm (Search)}
1.  Bắt đầu từ Root Page.
2.  So sánh khóa cần tìm với các khóa trong Node nội bộ để chọn đúng PageId của Node con.
3.  Lặp lại cho đến khi chạm tới Node lá.
4.  Tại Node lá, duyệt qua các Slot để lấy dữ liệu.
*   \textbf{Độ phức tạp:} $O(\log\_b n)$, với $b$ là bậc của cây (~600-1000 cho trang 16KB).

\subsubsection{Thuật toán Chèn và Tách Node (Insert & Split)}
Khi một Node (Lá hoặc Nội bộ) bị đầy và không thể chèn thêm bản ghi mới:
1.  \textbf{Split:} Tạo một Page mới cùng cấp.
2.  \textbf{Move:} Chuyển một nửa số lượng khóa sang Page mới.
3.  \textbf{Update Parent:} Chèn khóa phân tách vào Node cha để dẫn hướng đến Page mới vừa tạo.
4.  \textbf{Propagate:} Nếu Node cha cũng đầy, quá trình tách sẽ được lặp lại ngược lên trên cho tới tận Root. Nếu Root bị tách, chiều cao của cây sẽ tăng thêm 1 lớp.

\subsubsection{Sơ đồ Tách Node (Split)}
\begin{figure}[ht!]
\centering
\includegraphics[width=0.8\textwidth]{assets/btree\_split\_flow.png}
\caption{Cơ chế tách nút (Split) khi trang dữ liệu (Page) bị đầy (Overflow).}
\label{fig:7_4}
\end{figure}
\textit{Hình: Cơ chế Tách Nút (Split) khi Page bị đầy (Overflow)}

\subsubsection{Thuật toán Xóa và Kéo khít Node (Delete & Merge)}
Khi một bản ghi bị xóa khỏi Node lá, thuật toán sẽ kiểm tra tỷ lệ lấp đầy (Fill Factor). Nếu không gian trống giảm quá ngưỡng tối thiểu (Thường là $< 50\\%$ hoặc báo lỗi \textbf{Underflow}), KBMS thực hiện:
1.  \textbf{Redistribute (Mượn khóa từ Sibling):} Kiểm tra các Node ruột thịt liền kề (trái/phải). Nếu chúng có dư thừa không gian (khóa > 50\%), hệ thống sẽ "mượn" một số khóa tràn sang Node hiện tại để bù đắp, đồng thời cập nhật lại Key ranh giới ở Node Nội bộ (Internal Node) bên trên.
2.  \textbf{Merge (Gộp trang):} Nếu cả Node ruột thịt bên cạnh cũng đang thoi thóp (cận mức tối thiểu), hệ thống tiến hành nối hai Node đó làm một. Cập nhật con trỏ `NextPageId` để vá lại chuỗi Leaf.
3.  \textbf{Delete in Parent:} Khóa phân giới trên Node cha trỏ xuống Node vừa bị hủy sẽ bị gỡ bỏ. Chu trình giải phóng này có thể đệ quy ngược lên tầng Root. Trang dữ liệu thừa sẽ được giao nộp lại cho `Free Space Manager` để tái chế.

---

\subsection{Quản lý Index trong KBMS}

Trong KBMS, mỗi Concept mặc định sẽ có một "Clustered Index" dựa trên cột ID (thường là biến đầu tiên).
*   \textbf{Dữ liệu đi kèm Index:} Khác với "Secondary Index" (chỉ chứa RowId), Clustered Index của KBMS chứa \textbf{toàn bộ nội dung} của Tuple ngay tại các Node lá.
*   \textbf{Lợi ích:} Truy cập trực tiếp dữ liệu chỉ sau một lần duyệt cây, giảm thiểu số lượng Fetch Page từ Buffer Pool.

---

\subsection{Ví dụ minh họa}

Giả sử ta có Concept `Product` với `id` là khóa chính:
*   \textbf{Root:} Chứa khóa [100, 200, 300].
*   \textbf{Node con 1:} Chứa các sản phẩm có ID từ 1 đến 99.
*   \textbf{Node con 2:} Chứa các sản phẩm từ 100 đến 199.
*   ...
*   Khi cần tìm sản phẩm ID=150, hệ thống sẽ đi từ Root $\rightarrow$ chọn Node con 2 $\rightarrow$ đọc dữ liệu tại Node lá đó.

\section{Cấu trúc Trang vật lý (Physical Page Layout)}

KBMS V3 sử dụng mô hình phân trang (Paging) và cấu trúc \textbf{Slotted Page} để tối ưu hóa truy xuất ngẫu nhiên và quản lý không gian trống hiệu quả cho các bản ghi có độ dài biến thiên.

---

\subsection{Giải phẫu Trang (Page Anatomy)}

Một trang dữ liệu chuẩn (mặc định 8KB hoặc 16KB) được phân rã thành 4 thành phần chính: \textbf{Header}, \textbf{Slot Array}, \textbf{Free Space}, và \textbf{Tuples}.

\begin{figure}[ht!]
\centering
\includegraphics[width=0.8\textwidth]{assets/binary\_page\_layout.png}
\caption{Giải phẫu cấu trúc trang nhị phân Slotted-Page của KBMS.}
\label{fig:7_5}
\end{figure}

*   \textbf{Header (24 Bytes)}: Nằm ở đầu trang, chứa các metadata điều hướng.
*   \textbf{Slot Array}: Danh sách các ô nhớ trỏ tới dữ liệu, phát triển theo chiều \textbf{tiến} (Forward). 
*   \textbf{Free Space}: Vùng nhớ trống nằm giữa Slot Array và Tuples.
*   \textbf{Tuples}: Dữ liệu thực thể thực tế, được chèn theo chiều \textbf{lùi} (Backward) từ cuối trang lên trên.

---

\subsection{Ví dụ Tính toán (Case Study)}

Giả sử chúng ta có Concept `Person` với các thuộc tính: `ID (Int 4B)`, `Age (Int 4B)`, `Salary (Double 8B)`. Tổng kích thước một bản ghi (Tuple) là \textbf{16 bytes}.

\textbf{Thông số kỹ thuật:}
1. \textbf{Page Size}: 16,384 Bytes.
2. \textbf{Header}: 24 Bytes.
3. \textbf{Slot Size}: 8 Bytes (dùng để quản lý vị trí và độ dài của Tuple).
4. \textbf{Chi phí cho mỗi bản ghi}: 16B (Dữ liệu) + 8B (Slot) = \textbf{24 Bytes}.

\textbf{Tính toán số lượng bản ghi tối đa trên một trang:}
\begin{equation}
MaxRecords = \lfloor (16384 - 24) / (16 + 8) \rfloor = \lfloor 16360 / 24 \rfloor = \mathbf{681}
\end{equation}

\textbf{Phân bổ không gian:}
*   \textbf{Header}: 24 Bytes.
*   \textbf{Slot Array}: $681 \times 8 = 5448$ Bytes.
*   \textbf{Tuples Storage}: $681 \times 16 = 10896$ Bytes.
*   \textbf{Tổng cộng}: $24 + 5448 + 10896 = 16368$ Bytes.
*   \textbf{Free Space dư thừa}: $16384 - 16368 = \mathbf{16}$ Bytes.

---

\subsection{Chi tiết Header (24 Bytes Structure)}

Thông số này được đối soát trực tiếp với mã nguồn `SlottedPage.cs`:

| Byte Offset | Trường dữ liệu | Kiểu dữ liệu | Ý nghĩa |
| :--- | :--- | :--- | :--- |
| \textbf{0 - 3} | \textbf{PageId} | int32 | ID duy nhất của trang trong tệp dữ liệu. |
| \textbf{4 - 7} | \textbf{LSN} | int32 | Số tuần tự nhật ký (Log Sequence Number) cho WAL. |
| \textbf{8 - 11} | \textbf{PrevPageId} | int32 | Con trỏ tới trang trước trong chuỗi B+ Tree. |
| \textbf{12 - 15} | \textbf{NextPageId} | int32 | Con trỏ tới trang sau phục vụ quét tuần tự. |
| \textbf{16 - 19} | \textbf{FreeSpacePointer} | int32 | Offset đánh dấu ranh giới bắt đầu của Tuple cuối cùng. |
| \textbf{20 - 23} | \textbf{TupleCount} | int32 | Tổng số lượng bản ghi (Slots) hiện có. |

---

\subsection{Cơ chế Slotted Page (Slot Operations)}

Để quản lý các bản ghi có độ dài khác nhau mà không làm phân mảnh trang:
*   \textbf{Slot Array}: Mỗi slot chiếm \textbf{8 bytes}: [Offset (4B) | Length (4B)].
*   \textbf{Record ID (RID)}: Một RID cụ thể sẽ bao gồm `(PageId, SlotId)`.
*   \textbf{Xóa bản ghi (Delete)}: Thay vì dịch chuyển dữ liệu, KBMS chỉ cần đánh dấu Offset và Length của Slot đó về `0`. Không gian trống sẽ được thu hồi khi thực hiện lệnh `VACUUM`.

\subsection{Lợi ích kỹ thuật}

1.  \textbf{Direct I/O}: Kích thước trang khớp với block size của hệ điều hành, giảm thiểu "Write Amplification".
2.  \textbf{Binary Search}: Slot Array cho phép tìm kiếm nhị phân bản ghi ngay trong trang với tốc độ cực cao.
3.  \textbf{Persistence}: Cơ chế LSN trong Header kết nối trực tiếp với hệ thống WAL, đảm bảo tính bền vững (Durability) của dữ liệu.

---

> [!IMPORTANT]
> \textbf{Lưu ý}: Việc thay đổi cấu trúc Header từ 9B sang 24B trong phiên bản V3 là bước cải tiến quan trọng để hỗ trợ phục hồi dữ liệu và quản lý giao dịch tri thức phức tạp.

\section{Độ bền & Phục hồi sau sự cố (Durability & Recovery)}

Đảm bảo dữ liệu không bị mất mát khi hệ thống bị ngắt điện đột ngột là nhiệm vụ của `WalManager`. Tài liệu này chi tiết hóa cơ chế phục hồi của KBMS V3.

---

\subsection{Giao thức Write-Ahead Logging (WAL)}

KBMS tuân thủ nguyên tắc: \textbf{Ghi Nhật ký trước khi ghi Dữ liệu}.

-   \textbf{Log First}: Mọi thay đổi (Insert, Update, Delete) được mô tả thành một đối tượng `LogRecord` và ghi vào tệp `transactions.klf` trên đĩa.
-   \textbf{Atomic Write}: Một bản ghi WAL chỉ được coi là thành công khi toàn bộ các bytes đã được hệ điều hành Flush xuống đĩa cứng (xác nhận bởi `FileStream.Flush(true)`).
-   \textbf{Memory Sync}: Sau khi Log được ghi thành công, dữ liệu trong `BufferPool` mới được đánh dấu là `Dirty` và cập nhật.

---

\subsection{Cơ chế Điểm kiểm tra (Checkpointing)}

Để ngăn tệp WAL (`transactions.klf`) phình to vô hạn, KBMS thực hiện quy trình \textbf{Checkpoint}:

1.  \textbf{Dừng ghi tạm thời}: Tạm dừng các giao dịch mới.
2.  \textbf{Flush Pages}: Ghi tất cả các trang bị thay đổi (`Dirty Pages`) trong Buffer Pool xuống tệp dữ liệu chính (`objects.kdf`).
3.  \textbf{Trị tiêu Log}: Xóa hoặc làm trống tệp `.klf` vì dữ liệu đã nằm an toàn trên file chính.
4.  \textbf{Ghi nhận trạng thái}: Cập nhật chỉ số Log mới nhất vào `metadata.kmf`.

---

\subsection{Quy trình Phục hồi (Restart Recovery)}

Khi KBMS khởi động lại sau một vụ treo máy (Crash), nó thực hiện quy trình \textbf{Forward Recovery}:

![Sơ đồ Vòng đời Phục hồi Dữ liệu (Recovery Flow)](../assets/diagrams/recovery\_flow.png)

<details>
<summary>Cấu trúc Mermaid (Source)</summary>

\begin{verbatim}
sequenceDiagram
    participant DB as KBMS Engine
    participant WAL as transactions.klf
    participant DATA as objects.kdf

    DB->>WAL: Open()
    DB->>WAL: Read from Last Checkpoint
    loop Với mỗi LogRecord hợp lệ
        DB->>DATA: Replay(LSN, Operation)
    end
    DB->>DATA: Sync & Checksum
    DB->>DB: System Ready
\end{verbatim}
</details>

-   \textbf{Redo Phase}: Hệ thống đọc tệp `.klf` từ vị trí Checkpoint cuối cùng và thực hiện lại mọi thao tác đã ghi trong Log để khớp với trạng thái RAM trước khi Crash.
-   \textbf{LSN (Log Sequence Number)}: Mỗi trang dữ liệu đều lưu trữ một chỉ số LSN tại \textbf{Byte Offset 4-7} trong Header 24B. Khi thực hiện phục hồi, hệ thống chỉ Replay các LogRecord có LSN lớn hơn LSN hiện tại của trang đó trên đĩa.
-   \textbf{Consistency Check}: Nếu bản ghi Log cuối cùng bị lỗi (Partial write), hệ thống sẽ bỏ qua bản ghi đó và Rollback về trạng thái nhất quán gần nhất.

---

\subsection{Phân tích Hiệu năng}

Ghi WAL là thao tác \textbf{Sequential I/O} (Ghi tuần tự), do đó nó có tốc độ cực nhanh (thường là $O(1)$) so với việc ghi ngẫu nhiên (Random I/O) vào cây B+ Tree. Đây là lý do KBMS có throughput ghi cao ngay cả trên các ổ đĩa HDD truyền thống.

\section{Quản lý Đĩa & Mã hóa (Disk Management & Encryption)}

Tầng thấp nhất của Storage Engine chịu trách nhiệm giao tiếp trực tiếp với hệ điều hành và bảo vệ dữ liệu tĩnh (Data-at-Rest).

---

\subsection{Ánh xạ Trang vật lý (Physical Block Mapping)}

KBMS V3 không lưu trữ dữ liệu dưới dạng luồng liên tục (Stream) mà chia thành các khối cố định (Blocks).

- \textbf{Page ID to Offset}: Mọi yêu cầu đọc trang `PageId` được `DiskManager` chuyển đổi thành vị trí file bằng công thức:
  \begin{equation}
Offset = PageId \times 16416
\end{equation}
- \textbf{Random Access ($O(1)$)}: Sử dụng phương thức `FileStream.Seek` để nhảy trực tiếp đến khối dữ liệu cần thiết mà không phải quét toàn bộ tệp, đảm bảo hiệu năng ổn định ngay cả khi tệp CSDL lên đến hàng trăm GB.

---

\subsection{Mã hóa Lưu trữ (AES Encryption)}

Mọi dữ liệu trong KBMS V3 đều được mã hóa ở mức trang (Page-level encryption) trước khi ghi xuống đĩa cứng.

- \textbf{Thuật toán}: AES-256.
- \textbf{Quy trình ghi}: 
    1. Lấy dữ liệu 16KB từ Buffer Pool.
    2. Mã hóa dữ liệu bằng `MasterKey`.
    3. Thêm 32 bytes IV (Initialization Vector) và Padding.
    4. Ghi khối 16,416 bytes xuống đĩa.
- \textbf{Quy trình đọc}: Thực hiện ngược lại: Đọc khối 16,416 bytes -> Giải mã -> Đưa 16KB dữ liệu thô vào RAM.

---

\subsection{Đặc tả Khối dữ liệu đĩa (Disk Block Spec)}

Do có lớp mã hóa, kích thước thực tế trên đĩa của một trang lớn hơn kích thước logic trong RAM.

| Thành phần | Kích thước (Bytes) | Mô tả |
| :--- | :--- | :--- |
| \textbf{Data Payload (RAM)} | 16,384 | 16KB dữ liệu thật (đã bao gồm 24B Header). |
| \textbf{AES IV/Padding (Disk)} | 32 | Phần bù mã hóa và vector khởi tạo 256-bit. |
| \textbf{Tổng cộng (On-Disk)} | \textbf{16,416} | \textbf{Kích thước một khối vật lý thực tế}. |

---

\subsection{Tính Toàn vẹn & Hiệu năng}

- \textbf{Flush to Disk}: Tham số `flushToDisk: true` được sử dụng trong mọi lệnh `WritePage` để yêu cầu Hệ điều hành ghi trực tiếp xuống phiến đĩa (Physical Media), tránh mất dữ liệu do cache của OS khi mất điện.
- \textbf{Zero-fill Allocation}: Khi cấp phát trang mới, `DiskManager` thực hiện ghi đè các bytes `0x00` để đảm bảo không gian đĩa được giữ chỗ trước (Pre-allocation), tránh phân mảnh tệp.

\section{Định dạng Tệp tin (File Formats)}

Hệ thống KBMS quản lý tri thức bền vững thông qua các loại tệp tin chuyên biệt. Việc hiểu rõ vai trò của từng loại tệp giúp quản trị viên thực hiện công tác bảo trì và sao lưu hiệu quả.

---

\subsection{Phân loại tệp tin chính}

| Định dạng | Tên gọi | Vai trò kỹ thuật |
| :--- | :--- | :--- |
| \textbf{.dat} | \textbf{Data File} | Chứa dữ liệu thực tế (Instances) của các Concept. Được cấu trúc dưới dạng B+ Tree nhị phân. |
| \textbf{.kbf} | \textbf{Knowledge Base File} | Lưu trữ định nghĩa Schema, Rules, Equations và Relations (Metadata). |
| \textbf{.log} | \textbf{WAL Log File} | Nhật ký ghi trước mọi thay đổi I/O để phục hồi trạng thái khi gặp sự cố đột ngột. |
| \textbf{.kbql} | \textbf{Script File} | Tệp văn bản chứa chuỗi câu lệnh KBQL hỗ trợ chế độ chạy hàng loạt (Batch Execution). |
| \textbf{.ini} | \textbf{Config File} | Tệp cấu hình hệ thống (Cổng kết nối, Đường dẫn dữ liệu, Master Key). |

\subsection{Tương tác giữa các tệp tin}

Khi người dùng thực hiện lệnh `USE KB Project\_Alpha`, Server sẽ kích hoạt chuỗi hành động sau:
1.  \textbf{Nạp Metadata}: Đọc tệp `Project\_Alpha.kbf` để tái cấu trúc mô hình tri thức trong bộ nhớ RAM (Parser/Inference Engine).
2.  \textbf{Kết nối Data}: Trỏ `DiskManager` tới tệp `Project\_Alpha.dat` để sẵn sàng truy xuất các trang dữ liệu qua Buffer Pool.
3.  \textbf{Kích hoạt WAL}: Khởi tạo hoặc tiếp tục ghi vào `Project\_Alpha.log` để bảo vệ các giao dịch tri thức sắp tới.

\subsection{Bảo mật tệp tin (Encryption at Rest)}

Mọi tệp \textbf{.dat} và \textbf{.kbf} đều được mã hóa bằng thuật toán \textbf{AES-256}.
*   \textbf{Master Key}: Được lưu trữ trong tệp cấu hình `.ini` hoặc biến môi trường.
*   \textbf{Data Key}: Mỗi KB có thể có một khóa con (Sub-key) riêng để tăng cường tính bảo mật đa tầng.

---

> [!TIP]
> \textbf{Sao lưu (Backup)}: Một bản sao lưu KB đầy đủ phải bao gồm cả hai tệp \textbf{.dat} và \textbf{.kbf} tương ứng để đảm bảo tính toàn vẹn của cả Dữ liệu và Tri thức.

\section{Xác thực Lưu trữ (Storage Validation)}

Tầng lưu trữ là "pháo đài" cuối cùng bảo vệ dữ liệu, do đó nó được kiểm thử với cường độ cao nhất về mặt hiệu năng I/O và độ bền.

\subsection{Kiểm thử Đơn vị Trang (Page-level Testing)}

Tệp \textbf{`StorageV3Tests.cs`} thực hiện kiểm tra việc đọc/ghi các Slotted Pages 16KB.

*   \textbf{Tính nguyên tố}: Đảm bảo toàn bộ 16,384 bytes được ghi xuống đĩa trọn vẹn.
*   \textbf{Hiệu năng I/O}: Thời gian ghi Slotted Page trung bình là \textbf{23ms}, thời gian thu hồi trang (Eviction) là \textbf{46ms}.
*   \textbf{Mã hóa}: Kiểm tra tính năng AES-256 (Dữ liệu sau khi ghi xuống đĩa không thể đọc được bằng trình xem HEX nếu không có khóa).

\subsubsection{Minh chứng Mã nguồn (Storage I/O):}
\textbf{[Missing Image: code_test_storage_io.png]}\\ \textit{Hình 7.6: Kịch bản kiểm thử luồng I/O và Buffer Pool.}

\subsubsection{Minh chứng Lưu trữ (HEX View):}
\textbf{[Missing Image: terminal_test_storage_hex.png]}\\ \textit{Hình 7.7: Kết quả đọc trang nhị phân trực tiếp từ đĩa cứng.}

\subsection{Kiểm thử Chỉ mục B+ Tree}

Kiểm tra khả năng tìm kiếm và phân tách nút (Node Splitting) khi dữ liệu vượt ngưỡng.

| Chỉ số kiểm thử | Giá trị thực tế | Mục tiêu |
| :--- | :--- | :--- |
| \textbf{Số lượng bản ghi} | 1,000,000 | Thỏa mãn thực tế |
| \textbf{Độ cao của cây} | 3 - 4 tầng | Tối ưu truy xuất |
| \textbf{Thời gian tìm kiếm} | < 1ms | Hiệu năng cao |

\subsubsection{Minh chứng Mã nguồn (B+ Tree):}
\textbf{[Missing Image: code_test_index.png]}\\ \textit{Hình 7.8: Kịch bản kiểm thử hiệu năng chỉ mục B+ Tree.}

\subsubsection{Minh chứng Log (Split Node):}
\textbf{[Missing Image: result_test_index.png]}\\ \textit{Hình 7.9: Kết quả tìm kiếm và duyệt cây B+ Tree.}

\subsection{Kiểm thử Hồi phục (WAL Recovery)}

Sử dụng tệp \textbf{`WalManagerV3`} để mô phỏng sự cố:
1.  Bắt đầu giao dịch (Start Transaction).
2.  Ghi dữ liệu vào RAM (Dirty Pages).
3.  Mô phỏng sập nguồn (Kill Process).
4.  Khởi động lại và kiểm tra việc \textbf{Redo} từ nhật ký `.log`.

\subsubsection{Minh chứng Mã nguồn (WAL Recovery):}
\textbf{[Missing Image: code_test_recovery.png]}\\ \textit{Hình 7.10: Kịch bản kiểm thử khả năng phục hồi dữ liệu từ file WAL.}

\subsubsection{Minh chứng Kết quả (Redo Log):}
\textbf{[Missing Image: result_test_recovery.png]}\\ \textit{Hình 7.11: Bằng chứng phục hồi trạng thái Concept thành công sau khi giả lập mất điện.}

---

> [!IMPORTANT]
> Toàn bộ các thử nghiệm trên xác nhận rằng KBMS V3 đạt tiêu chuẩn ACID về độ bền vững và an toàn dữ liệu.

